% Generated draft: requires mathematical and source review.
\documentclass[11pt]{article}
\usepackage[margin=25mm]{geometry}
\usepackage{fontspec}
\setmainfont{Latin Modern Roman}
\usepackage{amsmath,amssymb,mathtools,unicode-math}
\setmathfont{Latin Modern Math}
\usepackage{enumitem,needspace,xurl,hyperref}
\setlength{\parindent}{0pt}
\setlength{\parskip}{5pt}
\setlength{\emergencystretch}{3em}
\newcommand{\N}{\mathbb{N}}
\newcommand{\Z}{\mathbb{Z}}
\newcommand{\Q}{\mathbb{Q}}
\newcommand{\R}{\mathbb{R}}
\newcommand{\C}{\mathbb{C}}
\newcommand{\F}{\mathbb{F}}
\newcommand{\A}{\mathbb{A}}
\newcommand{\PP}{\mathbb{P}}
\newcommand{\E}{\mathbb{E}}
\newcommand{\eps}{\varepsilon}
\newcommand{\dd}{\,\mathrm{d}}
\begin{document}
\section*{UnsolvedMath ID: 800095}
\textbf{Source problem code:} catalogue entry
\subsection{Definitions and mathematical statement}
Let \(T^2=\R^2/\Z^2\), and let \(p:E\to S\) be the torus bundle in the supplied question. The source denotes its monodromy by \(\rho\), but does not specify the hypotheses on \(S\), the bundle's structure group, or the target group of \(\rho\). A compact leaf of a foliation means a leaf that is compact in its intrinsic leaf topology.

\textbf{Question.} Under what hypotheses on the monodromy \(\rho\) does \(E\) admit a foliation \(\mathcal F\) with no compact leaves that is ``close to a fiber''? The no-compact-leaves requirement is
\[
\text{for every leaf }L\text{ of }\mathcal F,\qquad L\text{ is noncompact}.
\]
The quoted closeness condition is unresolved: the supplied excerpt specifies neither a topology nor a quantitative neighborhood, and does not say whether arbitrarily close foliations are required. It also does not explicitly specify the dimension or regularity of \(\mathcal F\). Consequently, the existence question cannot be given a fully quantified precise formulation from this excerpt alone.

\textbf{Supplied vector-field context.} For vector fields \(X,Y\) on a common manifold \(M\), commuting means \([X,Y]=0\), and a common zero is a point \(x\in M\) with \(X(x)=Y(x)=0\). The excerpt credits Elon Lima with a common-zero result for commuting vector fields in the trivial-representation case. The underlying manifold, regularity, and other hypotheses of that result are absent, so this attribution must not be read as an unrestricted common-zero theorem.

The excerpt also contains a first-person report of an example of a \(C^1\) commuting pair with no common zero, equivariant with respect to \(\rho\), when \(\rho(\pi_1(S))\) is generated by a single hyperbolic matrix. In the usual linear terminology, a hyperbolic matrix has no eigenvalue of modulus \(1\). The supplied excerpt does not identify the speaker, the vector fields' domain, the relevant group actions, or the precise equivariance condition. It describes the torus-bundle question as \emph{essentially equivalent} to the preceding problem; the missing setting prevents a precise equivalence statement.
\subsection{Short English statement}
Determine which monodromies of a torus bundle over \(S\) allow a foliation without compact leaves that is close to a fiber. The supplied excerpt leaves the assumptions on \(S\), foliation regularity and dimension, and meaning of closeness unspecified. It also reports related common-zero and equivariant vector-field results whose full hypotheses are missing.
\subsection{Sources}
\begin{itemize}
\item Local UnsolvedMath registry, permanent ID \(800095\), source code \texttt{catalogue entry}, title ``Foliations and Diffeomorphism Groups --- Question 36.'' The supplied registry attributes the record to Nariman and Yazdi, with proposed year \(2025\), and labels it open. This label is supplied source metadata, not an independently verified status assessment.
\item Nariman and Yazdi, \emph{Problem List on Foliations and Diffeomorphism Groups} (2025), Question 36. Supplied source URL: \url{https://arxiv.org/abs/2503.18274}. The record reports access on 2026-10-16 and extraction from arXiv LaTeX; these are supplied provenance statements, not a claim of independent source consultation.
\item The supplied excerpt credits Elon Lima with the trivial-representation common-zero result, without giving a bibliographic citation or its complete hypotheses. The first-person example and the claimed essential equivalence are preserved as reports from that excerpt.
\end{itemize}
\end{document}
